\section{Introduction}

\subsection{Project Background}

With the rapid development of financial markets, the number and complexity of
financial research reports are increasing. Analysts and investors need to
quickly extract key information from a large number of reports to make informed
investment decisions. However, traditional manual reading methods are
inefficient and difficult to handle massive amounts of information.

In recent years, the development of large language models (LLM) and
retrieval-augmented generation (RAG) technology has provided new solutions for
financial report analysis. By combining the capabilities of vector retrieval
and large language models, intelligent analysis and question-answering for
financial reports can be achieved.

\subsection{Project Objectives}

The core objective of the Insight: Financial Report Analysis Engine project is
to build a vertical domain RAG question-answering system that benchmarks the
NotebookLM experience, focusing on the analysis of long, information-dense
financial PDF reports, and providing page-level citation tracking.

Specific objectives include:

\begin{enumerate}
    \item Implement efficient parsing and processing of financial PDF reports
    \item Build high-quality vector indexes to support fast and accurate information
          retrieval
    \item Integrate local large language models to provide intelligent question-answering
          and analysis capabilities
    \item Implement page-level citation tracking to enhance the credibility of answers
    \item Establish a comprehensive evaluation system to ensure system performance
\end{enumerate}

\subsection{Technology Stack}

The project adopts the following technology stack:

\begin{itemize}
    \item \textbf{Frontend Framework}: Gradio
    \item \textbf{Backend Framework}: LlamaIndex v0.10+
    \item \textbf{Vector Database}: ChromaDB
    \item \textbf{Embedding Model}: BAAI/bge-m3
    \item \textbf{Large Language Model}: Qwen2.5-7B/72B (locally deployed via Ollama)
    \item \textbf{PDF Parsing}: PyMuPDF
    \item \textbf{Evaluation Framework}: RAGAS
\end{itemize}

\subsection{Project Structure}

The project adopts a modular design, mainly consisting of the following
directory structure:

\begin{itemize}
    \item `src/ingest/`: Data ingestion layer, responsible for PDF parsing and chunking
    \item `src/retrieval/`: Retrieval layer, implementing vector retrieval and hybrid retrieval
    \item `src/generation/`: Generation layer, responsible for building the dialogue engine and generating answers
    \item `src/evaluation/`: Evaluation layer, implementing RAGAS evaluation
    \item `src/ui/`: Interaction layer, building Gradio frontend
    \item `src/utils/`: Utility modules, providing configuration management and other functions
    \item `configs/`: Configuration file directory
    \item `data/`: Data file directory
    \item `docs/`: Documentation directory
\end{itemize}
